\subsection{Counterexamples to stronger readings}
\hand \textbf{No universal flood price.} Book many registrations without
folding them. No $R$ donation occurs. Even eventual balanced self-rejection
can make the tip internal, leaving only fees and opportunity cost. Successful
junk registrations do pay $R$ but still consume state-output turns.

\hand \textbf{Protection is not priority.} Two authorized registrations for
the same unknown key can be pending. The first fold makes it active; the
second is now mismatching and may be rejected before its time expiry. Its
refund floor survives, not a right to the key or the full tip.

\hand \textbf{Noninterference needs disjoint inputs.} If a proposed fold and
a rejection spend the same pending registration, rejection removes that
input. Preserving the state output alone cannot preserve the whole fold.
Conversely repeated folds can invalidate a rejection's state reference before
it lands, without violating the rejection theorem.

\hand \textbf{No positive wage.} Choose fold tip $t=100{,}000$, fee
$f=170{,}000$ and ordinary funding $I\ge170{,}000$. The allocation equation
gives net ordinary change $-70{,}000$. Fold acceptance requires funding but
does not impose the availability theorem's extra $f\le t$ hypothesis.

\hand \textbf{No ownership or squatting defence.} An application that
authorizes every request also authorizes a stranger's registration of a
desirable key. Paying $R$ does not prove a right to its name. A tag's unique
custody is an asset invariant, not an ownership adjudication.

\missing \textbf{Full refund is not fee-paying retractability.} A single
$3{,}000{,}000$ input cannot return that full amount and also pay $170{,}000$.
The current accepted-retraction theorem and this conservation counterexample
coexist because its judgement omits balance and excludes a funding input.
Changing the model or claiming real-ledger retraction requires a separate
ruling and new evidence; the floor is not silently weakened here.
